\section{定位：防御型静默吸纳，而非高增长反转}

沪农商行是资金行为最“安静”的标的：连续10日净流入1.42亿元，同期股价反而下跌3.04\%；7月10日收盘7.66元，位于过去一年价格区间约8\%。融资余额仅1.31亿元，占流通市值约0.18\%，并且当日融资净偿还。它更像保险、国资和收益型资金承接的低波动资产，而不是游资推动的主题行情。

\section{基本面：改善存在，但弹性有限}

\begin{exhibitbox}[Exhibit 8：沪农商行2026Q1关键指标]
\begin{tabularx}{\textwidth}{L{3.4cm}R{2.5cm}X}
\toprule
\textbf{指标} & \textbf{数据} & \textbf{解读} \\
\midrule
营业收入 & 66.41亿元/+1.23\% & 从2025全年负增长转正 \\
归母净利润 & 35.90亿元/+0.73\% & 维持正增长但弹性弱 \\
净利息收入 & 同比+2.3\% & 息差压力收窄 \\
手续费及佣金净收入 & 同比+16.9\% & 保险、理财代销改善 \\
不良贷款率 & 0.96\% & 保持稳定 \\
拨备覆盖率 & 303.07\% & 仍高，但较年初下降 \\
2025分红率 & 34.07\% & 收益型资金的重要锚 \\
\bottomrule
\end{tabularx}
\sourcenote{公司2026Q1报告；国泰海通、国信、华创公开研究。}
\end{exhibitbox}

沪农商行的问题不是资产质量，而是盈利增长。国信预测2026E归母净利润约124亿元、EPS 1.29元，利润增速不足1\%；国泰海通预测也只有约1.1\%。因此，估值修复主要依靠PB、股息率和息差企稳，而不是利润高速增长。

\section{持股结构与风险边界}

2026Q1前十大股东以上海国资、央企、保险资金和长期法人为主，前十大没有沪股通。公开F10口径显示香港中央结算后续进入主要流通股东，但披露时点和口径不同，因此本报告不把北向作为核心结论。真正可确认的是：持有人结构偏长期、融资盘很低、连续大单流入尚未转化成价格趋势。

\section{估值与动作}

国信预测2026E BVPS 14.11元，国泰海通预测14.31元。AStock取基准0.70倍PB，对应9.88元；叠加8.50元市场锚和国泰海通10.73元外部目标，最终目标9.69元，相对7.66元空间约26.5\%。

\begin{itemize}
  \item \textbf{配置区}：7.55--7.80元，适合收益型资金。
  \item \textbf{确认区}：8.10--8.50元，需下一季营收保持正增长。
  \item \textbf{失效条件}：不良率升破1.05\%、拨备覆盖率跌破280\%、营收重新负增长或分红率下调。
\end{itemize}
